\documentclass[11pt]{article}
\usepackage[margin=1in]{geometry}
\usepackage{amsmath,amssymb,amsthm}
\usepackage[colorlinks=true,urlcolor=blue]{hyperref}
\newtheorem{theorem}{Theorem}
\newtheorem{lemma}[theorem]{Lemma}
\newtheorem{proposition}[theorem]{Proposition}
\newtheorem{remark}[theorem]{Remark}
\newcommand{\E}{\mathbb E}
\newcommand{\Prob}{\mathbb P}
\newcommand{\one}{\mathbf 1}
\title{An adaptive cross-matching defeats planted-kernel simplicity:\\
a clean short-cycle proof}
\author{Independent adversarial boundary audit}
\date{November 15, 2026}
\begin{document}
\maketitle

\paragraph{Scope.}
This note proves the requested boundary between an independent cross
graph and a cross graph selected after seeing the internal graphs.
The growing-cycle calculation below is self-contained in the cubic
configuration model. The only imported probability fact is the standard
positive simplicity probability, here cited precisely to
Janson~\cite[Corollary 1.5]{janson}.
No historical novelty claim is made.

\begin{lemma}[A clean cycle with high probability]\label{lem:cycle}
Let $N$ be even, and let $H_N$ be a uniform simple cubic graph on $N$
labelled vertices. Put $\ell=\lfloor\log\log N\rfloor$, with natural
logarithms. With probability tending to one, $H_N$ contains a cycle of
length exactly $\ell$ that has no chord and whose $\ell$ third neighbors
are outside the cycle and pairwise distinct.
\end{lemma}

\begin{proof}
We first work with the configuration model: three labelled half-edges
at each vertex, paired uniformly among the $3N$ half-edges. For a
positive integer $h$ define
\[
 D_h=\prod_{j=0}^{h-1}(3N-2j-1).
\]
Any specified $h$ disjoint half-edge pairs occur with probability $D_h^{-1}$.
Let $X$ count unoriented half-edge cycles of length $\ell$ with distinct
vertices. There are $(N)_\ell6^\ell/(2\ell)$ candidates, because at each
vertex its ordered incoming and outgoing half-edges have six choices.
Consequently
\begin{equation}\label{eq:mean}
 \mu:=\E X=\frac{(N)_\ell6^\ell}{2\ell D_\ell}
       =\frac{2^\ell}{2\ell}
                       \bigl(1+O(\ell^2/N)\bigr)\longrightarrow\infty.
\end{equation}

We bound the second factorial moment, retaining the dependence on
$\ell$. Ordered pairs of vertex-disjoint cycles contribute exactly
\[
 \frac{(N)_{2\ell}6^{2\ell}}{(2\ell)^2D_{2\ell}}
                =\mu^2\bigl(1+O(\ell^2/N)\bigr).
\]
For distinct cycles sharing a vertex, their union is connected and
contains two distinct cycles, so if the union has $v$ vertices and $e$
edges, then $e\ge v+1$. This also holds for half-edge multigraphs;
incompatible prescribed pairs have zero probability.

Here is a deliberately crude uniform bound on all such overlapping
pairs. Present both cycles by ordered vertex sequences. An identification
pattern chooses which positions in the second sequence coincide with
positions of the first, and which first positions they use. The number
of patterns is at most
\[
              \sum_{s=0}^{\ell}\binom\ell s\ell^s=(1+\ell)^\ell.
\]
For a pattern with $v$ union vertices there are at most $N^v$ vertex
labellings and at most $6^{2\ell}$ choices of cycle half-edges. Since
$e\le2\ell$ and eventually $N\ge2\ell$, each compatible distinct
overlapping pair occurs with probability
\[
                   D_e^{-1}\le N^{-e}\le N^{-v-1}.
\]
Ignoring the rotational/orientation divisions only increases this upper
bound. Thus the total overlapping-pair contribution is at most
\begin{equation}\label{eq:overlap}
           \frac{36^\ell(1+\ell)^\ell}{N}=N^{-1+o(1)}.
\end{equation}
Diagonal pairs contribute $\mu$ to $\E X^2$. Therefore
\[
 \frac{\operatorname{Var}X}{\mu^2}
 \le\frac1\mu+O(\ell^2/N)
       +\frac{36^\ell(1+\ell)^\ell}{N\mu^2}=o(1).
\]
Chebyshev's inequality gives $\Prob(X=0)=o(1)$.

We next show that with high probability every such cycle is clean.
Condition on one specified candidate cycle occurring. Its vertices have
one unused half-edge apiece, and the remaining pairing is uniform on
$R=3N-2\ell$ half-edges. A chord requires a pair among the $\ell$
unused cycle half-edges, with probability at most
\[
                         \frac{\binom\ell2}{R-1}.
\]
Two third neighbors coincide outside the cycle only if two distinct
cycle half-edges are paired to two distinct half-edges at the same
outside vertex. A union bound gives probability at most
\[
             \frac{6\binom\ell2(N-\ell)}{(R-1)(R-3)}.
\]
The sum is $O(\ell^2/N)$, uniformly over the specified cycle. If $Y$
counts bad cycles, then
\[
       \E Y\le O(\ell^2/N)\mu
                  =O(\ell2^\ell/N)=o(1).
\]
Markov therefore gives $\Prob(Y>0)=o(1)$. Combining with $X>0$ proves
the clean-cycle assertion for configurations.

Conditioning the configuration on simplicity yields the uniform simple
cubic graph. By Janson~\cite[Corollary 1.5]{janson}, the simplicity
probability tends to $e^{-2}>0$ (its parameter there equals one in the
cubic case). Dividing the just-proved failure probability by this
positive conditioning probability transfers the claim to $H_N$.
\end{proof}

\begin{remark}
The same estimates work for any integer sequence $\ell\to\infty$ with
$\ell\log\ell=o(\log N)$. No fixed-length Poisson approximation has
been applied outside its proved range; the second-moment calculation
explicitly controls the growing length.
\end{remark}

\begin{proposition}[Deterministic adaptive matching construction]
\label{prop:matching}
Let $C,D$ be cubic simple graphs on two labelled blocks $V_+,V_-$ of
equal size $N$. Suppose each contains a clean cycle of the same length
$\ell$, with $2\ell<N$. There exists a simple perfect cross-matching
$B$ such that the resulting graph has two different balanced
$(3,1)$-equitable partitions and its adjacency matrix has eigenvalue two
of multiplicity at least two.
\end{proposition}

\begin{proof}
Let $S\subset V_+$ and $T\subset V_-$ be the cycles, and let
$U\subset V_+\setminus S$ and $W\subset V_-\setminus T$ be their third
neighbor sets. Cleanliness gives $|S|=|T|=|U|=|W|=\ell$, with each
vertex of $U$ adjacent to exactly one vertex of $S$, and likewise for
$W,T$.

Choose arbitrary bijections matching $S$ to $W$ and $U$ to $T$.
Complete this to a perfect cross-matching by any bijection from
$V_+\setminus(S\cup U)$ to $V_-\setminus(T\cup W)$.
Every vertex then has three internal neighbors and one cross neighbor,
so the original block partition is balanced $(3,1)$-equitable.

Flip the membership of exactly $S\cup T$. A cycle vertex has two
internal neighbors that flip with it, one internal neighbor that does
not, and a cross-matching neighbor in the unflipped boundary set on the
other original block. It consequently has three same-side and one
opposite-side neighbors in the new partition. A boundary vertex has
one internal neighbor that flips, two that do not, and its matching
neighbor on the other cycle, which flips; it too has three same-side
and one opposite-side neighbors. Every remaining vertex has no internal
neighbor on either cycle and its matching neighbor also remains
unflipped, so its three-to-one degree pattern is preserved. Both sides
still have $N$ vertices, because $|S|=|T|$.

Equivalently, the vector $v$ equal to $+1$ on $S$, $-1$ on $T$, and
zero elsewhere satisfies $Av=2v$: cycle vertices see the two cycle
neighbors, and at boundary vertices the cycle and cross-matching
contributions cancel. Let $\sigma$ be $+1$ on $V_+$ and $-1$ on
$V_-$. Then $A\sigma=2\sigma$, while $v$ and $\sigma$ are linearly
independent since $0<2\ell<2N$. The alternative planted sign vector is
exactly $\sigma-2v$, proving both conclusions.
\end{proof}

\begin{theorem}[Failure under a cross graph selected adaptively]
Let $C,D$ be independent uniform simple cubic graphs on even block size
$N$. With probability tending to one, there exists a perfect cross
matching $B=B(C,D)$ for which
\[
 \operatorname{nullity}_{\mathbb R}
 \left(\begin{pmatrix}C&B\\B^T&D\end{pmatrix}-2I\right)\ge2.
\]
It can be chosen to produce a second admissible partition differing
from the original on exactly $2\lfloor\log\log N\rfloor$ vertices.
\end{theorem}
\begin{proof}
Apply Lemma~\ref{lem:cycle} to both internal graphs, then
Proposition~\ref{prop:matching}. Lexicographically choose the cycles
and bijections when they exist to obtain a definite selection function;
choose any matching on the vanishing-probability failure event.
\end{proof}

\paragraph{Quantifier boundary.}
This is consistent with a theorem saying
\[
 \sup_{B\text{ deterministic}}\Prob_{C,D}(\text{failure for }B)=o(1).
\]
The theorem here instead has the quantifiers
\[
 \Prob_{C,D}(\exists B:\text{failure for }B)=1-o(1).
\]
The adaptive construction depends on the actual internal graphs and
does not challenge a result for an independent cross graph. It exhibits
an explicit localized additional eigenvector, rather than inferring
kernel multiplicity merely from an approximate spectral argument.

\begin{thebibliography}{9}
\bibitem{janson}
S.~Janson, \emph{The probability that a random multigraph is simple},
arXiv:math/0609802v1, 2006, Corollary~1.5.
\url{https://arxiv.org/pdf/math/0609802v1}.
The only imported facts are positive fixed-degree simplicity probability
and the conditional uniform-simple model, stated in its introduction.
\end{thebibliography}
\end{document}
